\section{Procedure di registrazione dei test}

\subsection{Struttura di un caso di test}
% TODO: campi che descrivono ogni caso di test: numero identificativo, descrizione, precondizioni, valori di input, output attesi, postcondizioni attese; all'esecuzione si aggiungono output e postcondizioni riscontrati ed esito.

\subsection{Registrazione dei risultati}
% TODO: dove si registrano i test (es. GitHub issues, foglio condiviso), come si segnalano i malfunzionamenti e come si tiene traccia delle correzioni.

\subsection{Casi di test}
% TODO: elenco dei casi di test. Esito positivo = malfunzionamento rilevato.

\begingroup
\SetTblrTemplate{contfoot}{empty}
\SetTblrTemplate{conthead}{empty}
\DefTblrTemplate{conthead-text}{nocont}{}
\DefTblrTemplate{contfoot-text}{nocont}{}
\SetTblrTemplate{conthead-text}{nocont}
\SetTblrTemplate{contfoot-text}{nocont}
\begin{longtblr}[entry=none,label=none]{width=\linewidth,
             colspec={X[0.5,l]X[l]X[l]X[l]X[l]X[l]X[0.6,l]X[0.7,l]}, rowhead=1,
             cells={font=\footnotesize},
             row{1}={bg=arancio}, hline{2-Z}={wd=0.5pt, fg=arancio}}
TC & Precond. & Input & Output attesi & Output riscontrato & PostCond. & Esito & Priorità \\
 & & & & & & & \\
\end{longtblr}
\endgroup
